\documentclass[10pt,a4paper]{amsart}
\usepackage{amsmath,amssymb,amsthm,graphicx}
% Retained preamble for the Stacks Project chapter snapshot in this directory.
% The chapter's own preamble.tex is a copy of the shared Stacks Project preamble; the block
% below is sliced MECHANICALLY, line for line, from it. Nothing was added, removed or
% redefined; the wrapper class, the xr-hyper/hyperref wiring and the externaldocument list
% are omitted because the assembler supplies the document class.
%
% lines taken from preamble.tex: 3, 5, 9, 11, 12, 156, 157, 158, ... (51 lines total)
%
% package coverage restored from the chapter's own preamble (2026-10-13): xy, xy-2cell

\usepackage[all]{xy}
\xyoption{2cell}
\UseAllTwocells

\newenvironment{reference}{\comment}{\endcomment}
\newenvironment{slogan}{\comment}{\endcomment}
\newenvironment{history}{\comment}{\endcomment}
\theoremstyle{plain}
\newtheorem{theorem}[subsection]{Theorem}
\newtheorem{proposition}[subsection]{Proposition}
\newtheorem{lemma}[subsection]{Lemma}
\theoremstyle{definition}
\newtheorem{definition}[subsection]{Definition}
\newtheorem{example}[subsection]{Example}
\newtheorem{exercise}[subsection]{Exercise}
\newtheorem{situation}[subsection]{Situation}
\theoremstyle{remark}
\newtheorem{remark}[subsection]{Remark}
\newtheorem{remarks}[subsection]{Remarks}
\numberwithin{equation}{subsection}
\def\lim{\mathop{\mathrm{lim}}\nolimits}
\def\colim{\mathop{\mathrm{colim}}\nolimits}
\def\Spec{\mathop{\mathrm{Spec}}}
\def\Hom{\mathop{\mathrm{Hom}}\nolimits}
\def\Ext{\mathop{\mathrm{Ext}}\nolimits}
\def\SheafHom{\mathop{\mathcal{H}\!\mathit{om}}\nolimits}
\def\SheafExt{\mathop{\mathcal{E}\!\mathit{xt}}\nolimits}
\def\Sch{\mathit{Sch}}
\def\Mor{\mathop{\mathrm{Mor}}\nolimits}
\def\Ob{\mathop{\mathrm{Ob}}\nolimits}
\def\Sh{\mathop{\mathit{Sh}}\nolimits}
\def\NL{\mathop{N\!L}\nolimits}
\def\CH{\mathop{\mathrm{CH}}\nolimits}
\def\proetale{{pro\text{-}\acute{e}tale}}
\def\etale{{\acute{e}tale}}
\def\QCoh{\mathit{QCoh}}
\def\Ker{\mathop{\mathrm{Ker}}}
\def\Im{\mathop{\mathrm{Im}}}
\def\Coker{\mathop{\mathrm{Coker}}}
\def\Coim{\mathop{\mathrm{Coim}}}
\DeclareMathSymbol{\boxtimes}{\mathbin}{AMSa}{"02}
\def\QCohstack{\mathcal{QC}\!\mathit{oh}}
\def\Cohstack{\mathcal{C}\!\mathit{oh}}
\def\Spacesstack{\mathcal{S}\!\mathit{paces}}
\def\Quotfunctor{\mathrm{Quot}}
\def\Hilbfunctor{\mathrm{Hilb}}
\def\Curvesstack{\mathcal{C}\!\mathit{urves}}
\def\Polarizedstack{\mathcal{P}\!\mathit{olarized}}
\def\Complexesstack{\mathcal{C}\!\mathit{omplexes}}
\def\Pic{\mathop{\mathrm{Pic}}\nolimits}
\def\Picardstack{\mathcal{P}\!\mathit{ic}}
\def\Picardfunctor{\mathrm{Pic}}
\def\Deformationcategory{\mathcal{D}\!\mathit{ef}}
\usepackage{cleveref}
\crefname{originallemma}{original lemma}{original lemmas}
\Crefname{originallemma}{Original lemma}{Original lemmas}
\makeatletter
\newlabel{descent-lemma-compare-sites}{{070S}{source}{}{sourcepointer-1}{}}
\newlabel{descent-lemma-compare-sites@cref}{{[originallemma][0][]070S}{[1][0][]source}{}{}{}}
\makeatother
\crefname{originallemma}{original lemma}{original lemmas}
\Crefname{originallemma}{Original lemma}{Original lemmas}
\makeatletter
\newlabel{descent-proposition-equivalence-quasi-coherent}{{03DX}{source}{}{sourcepointer-2}{}}
\newlabel{descent-proposition-equivalence-quasi-coherent@cref}{{[originallemma][0][]03DX}{[1][0][]source}{}{}{}}
\makeatother
\crefname{originallemma}{original lemma}{original lemmas}
\Crefname{originallemma}{Original lemma}{Original lemmas}
\makeatletter
\newlabel{descent-remark-change-topologies-ringed-sites}{{03FH}{source}{}{sourcepointer-3}{}}
\newlabel{descent-remark-change-topologies-ringed-sites@cref}{{[originallemma][0][]03FH}{[1][0][]source}{}{}{}}
\makeatother
\crefname{originallemma}{original lemma}{original lemmas}
\Crefname{originallemma}{Original lemma}{Original lemmas}
\makeatletter
\newlabel{descent-section-quasi-coherent-sheaves}{{03DR}{source}{}{sourcepointer-4}{}}
\newlabel{descent-section-quasi-coherent-sheaves@cref}{{[originallemma][0][]03DR}{[1][0][]source}{}{}{}}
\makeatother
\crefname{originallemma}{original lemma}{original lemmas}
\Crefname{originallemma}{Original lemma}{Original lemmas}
\makeatletter
\newlabel{equation-restrict}{{equation-restrict}{source}{}{sourcepointer-5}{}}
\newlabel{equation-restrict@cref}{{[originallemma][0][]equation-restrict}{[1][0][]source}{}{}{}}
\makeatother
\crefname{originallemma}{original lemma}{original lemmas}
\Crefname{originallemma}{Original lemma}{Original lemmas}
\makeatletter
\newlabel{lemma-characterize-quasi-coherent-bis}{{lemma-characterize-quasi-coherent-bis}{source}{}{sourcepointer-6}{}}
\newlabel{lemma-characterize-quasi-coherent-bis@cref}{{[originallemma][0][]lemma-characterize-quasi-coherent-bis}{[1][0][]source}{}{}{}}
\makeatother
\crefname{originallemma}{original lemma}{original lemmas}
\Crefname{originallemma}{Original lemma}{Original lemmas}
\makeatletter
\newlabel{lemma-compare}{{lemma-compare}{source}{}{sourcepointer-7}{}}
\newlabel{lemma-compare@cref}{{[originallemma][0][]lemma-compare}{[1][0][]source}{}{}{}}
\makeatother
\crefname{originallemma}{original lemma}{original lemmas}
\Crefname{originallemma}{Original lemma}{Original lemmas}
\makeatletter
\newlabel{lemma-compare-locally-quasi-coherent}{{lemma-compare-locally-quasi-coherent}{source}{}{sourcepointer-8}{}}
\newlabel{lemma-compare-locally-quasi-coherent@cref}{{[originallemma][0][]lemma-compare-locally-quasi-coherent}{[1][0][]source}{}{}{}}
\makeatother
\crefname{originallemma}{original lemma}{original lemmas}
\Crefname{originallemma}{Original lemma}{Original lemmas}
\makeatletter
\newlabel{lemma-compare-morphism}{{lemma-compare-morphism}{source}{}{sourcepointer-9}{}}
\newlabel{lemma-compare-morphism@cref}{{[originallemma][0][]lemma-compare-morphism}{[1][0][]source}{}{}{}}
\makeatother
\crefname{originallemma}{original lemma}{original lemmas}
\Crefname{originallemma}{Original lemma}{Original lemmas}
\makeatletter
\newlabel{remark-compare}{{remark-compare}{source}{}{sourcepointer-10}{}}
\newlabel{remark-compare@cref}{{[originallemma][0][]remark-compare}{[1][0][]source}{}{}{}}
\makeatother
\crefname{originallemma}{original lemma}{original lemmas}
\Crefname{originallemma}{Original lemma}{Original lemmas}
\makeatletter
\newlabel{sites-modules-lemma-local-pullback}{{03DO}{source}{}{sourcepointer-11}{}}
\newlabel{sites-modules-lemma-local-pullback@cref}{{[originallemma][0][]03DO}{[1][0][]source}{}{}{}}
\makeatother
\title{Stacks, Sheaves on Algebraic Stacks: lemma-compare-quasi-coherent}
\author{The Stacks Project Authors}
\date{stacks-project:stacks-sheaves.tex:local-snapshot-20261012; GNU FDL 1.2 or later}
\begin{document}
\maketitle
\noindent\textit{Editorial note.} Statement span is lines 2262-2274 of sources/stacks-algebraic-stacks-sheaves/pinned-original-source.tex; proof span is lines 2276-2339, the complete author proof (62 body lines).\par\medskip
\noindent\textit{Editorial note.} Navigation stubs: 11 cited-but-not-included labels are supplied as non-load-bearing pointers with their pinned tags where the snapshot tag index has them.\par\medskip
\noindent\textit{Editorial note.} extra\_packages is empty on purpose: the retained preamble already loads xy with options and a second declaration clashes in the assembler wrapper.\par\medskip
\noindent\textit{Editorial note.} Text is retained verbatim, including the author's own typos and placeholders.\par\medskip
\noindent\textit{Editorial note.} Queue index 5 of queue\_algebraic-stacks-sheaves.json (pb=3017, stmt 2262-2274, proof 2276-2339).\par\medskip
\noindent\textit{Allowed standard prerequisites (editorial).}\begin{itemize}
\item Explicit prerequisite boundary. The cited inputs are prerequisites; the argument in the retained span is not replaced by any cited result.
\end{itemize}\medskip
\ifdefined\hypertarget\hypertarget{sourcepointer-1}{}\fi
\noindent\textit{Source reference.} Original lemma 070S: sources/stacks-algebraic-stacks-sheaves/pinned-original-source.tex, source lines 2339--2339. Cited by the retained statement or proof but not part of the collected spans; supplied only as a navigation pointer. The argument does not depend on its content.\par\medskip
\ifdefined\hypertarget\hypertarget{sourcepointer-2}{}\fi
\noindent\textit{Source reference.} Original lemma 03DX: sources/stacks-algebraic-stacks-sheaves/pinned-original-source.tex, source lines 2339--2339. Cited by the retained statement or proof but not part of the collected spans; supplied only as a navigation pointer. The argument does not depend on its content.\par\medskip
\ifdefined\hypertarget\hypertarget{sourcepointer-3}{}\fi
\noindent\textit{Source reference.} Original lemma 03FH: sources/stacks-algebraic-stacks-sheaves/pinned-original-source.tex, source lines 2339--2339. Cited by the retained statement or proof but not part of the collected spans; supplied only as a navigation pointer. The argument does not depend on its content.\par\medskip
\ifdefined\hypertarget\hypertarget{sourcepointer-4}{}\fi
\noindent\textit{Source reference.} Original lemma 03DR: sources/stacks-algebraic-stacks-sheaves/pinned-original-source.tex, source lines 2339--2339. Cited by the retained statement or proof but not part of the collected spans; supplied only as a navigation pointer. The argument does not depend on its content.\par\medskip
\ifdefined\hypertarget\hypertarget{sourcepointer-5}{}\fi
\noindent\textit{Source reference.} Original lemma equation-restrict: sources/stacks-algebraic-stacks-sheaves/pinned-original-source.tex, source lines 1547--1549. Cited by the retained statement or proof but not part of the collected spans; supplied only as a navigation pointer. The argument does not depend on its content.\par\medskip
\ifdefined\hypertarget\hypertarget{sourcepointer-6}{}\fi
\noindent\textit{Source reference.} Original lemma lemma-characterize-quasi-coherent-bis: sources/stacks-algebraic-stacks-sheaves/pinned-original-source.tex, source lines 1819--1836. Cited by the retained statement or proof but not part of the collected spans; supplied only as a navigation pointer. The argument does not depend on its content.\par\medskip
\ifdefined\hypertarget\hypertarget{sourcepointer-7}{}\fi
\noindent\textit{Source reference.} Original lemma lemma-compare: sources/stacks-algebraic-stacks-sheaves/pinned-original-source.tex, source lines 1430--1452. Cited by the retained statement or proof but not part of the collected spans; supplied only as a navigation pointer. The argument does not depend on its content.\par\medskip
\ifdefined\hypertarget\hypertarget{sourcepointer-8}{}\fi
\noindent\textit{Source reference.} Original lemma lemma-compare-locally-quasi-coherent: sources/stacks-algebraic-stacks-sheaves/pinned-original-source.tex, source lines 2247--2253. Cited by the retained statement or proof but not part of the collected spans; supplied only as a navigation pointer. The argument does not depend on its content.\par\medskip
\ifdefined\hypertarget\hypertarget{sourcepointer-9}{}\fi
\noindent\textit{Source reference.} Original lemma lemma-compare-morphism: sources/stacks-algebraic-stacks-sheaves/pinned-original-source.tex, source lines 1553--1570. Cited by the retained statement or proof but not part of the collected spans; supplied only as a navigation pointer. The argument does not depend on its content.\par\medskip
\ifdefined\hypertarget\hypertarget{sourcepointer-10}{}\fi
\noindent\textit{Source reference.} Original lemma remark-compare: sources/stacks-algebraic-stacks-sheaves/pinned-original-source.tex, source lines 1508--1529. Cited by the retained statement or proof but not part of the collected spans; supplied only as a navigation pointer. The argument does not depend on its content.\par\medskip
\ifdefined\hypertarget\hypertarget{sourcepointer-11}{}\fi
\noindent\textit{Source reference.} Original lemma 03DO: sources/stacks-algebraic-stacks-sheaves/pinned-original-source.tex, source lines 2339--2339. Cited by the retained statement or proof but not part of the collected spans; supplied only as a navigation pointer. The argument does not depend on its content.\par\medskip
\begin{lemma}
\label{lemma-compare-quasi-coherent}
Let $S$ be a scheme. Let $\mathcal{X} \to (\Sch/S)_{fppf}$ be a category
fibred in groupoids which is representable by an algebraic space $F$.
The functor (\ref{equation-restrict}) defines an equivalence
$$
\QCoh(\mathcal{O}_\mathcal{X}) \to \QCoh(\mathcal{O}_F),\quad
\mathcal{F} \longmapsto \mathcal{F}|_{F_\etale}
$$
with quasi-inverse given by $\mathcal{G} \mapsto \pi_F^*\mathcal{G}$.
This equivalence is compatible with pullback for morphisms between
categories fibred in groupoids representable by algebraic spaces.
\end{lemma}

\begin{proof}
By Lemma \ref{lemma-characterize-quasi-coherent-bis} we may work with
the \'etale topology. We will use the notation and results of
Lemma \ref{lemma-compare} without further mention.
Recall that the restriction functor
$\textit{Mod}(\mathcal{X}_\etale, \mathcal{O}_\mathcal{X}) \to
\textit{Mod}(F_\etale, \mathcal{O}_F)$,
$\mathcal{F} \mapsto \mathcal{F}|_{F_\etale}$ is given by $i_F^*$. By
Lemma \ref{lemma-compare-locally-quasi-coherent} or by
Modules on Sites, Lemma \ref{sites-modules-lemma-local-pullback}
we see that $\mathcal{F}|_{F_\etale}$ is quasi-coherent if
$\mathcal{F}$ is quasi-coherent.
Hence we get a functor as indicated in the statement of the
lemma and we get a functor $\pi_F^*$ in the opposite direction.
Since $\pi_F \circ i_F = \text{id}$ we see that
$i_F^*\pi_F^*\mathcal{G} = \mathcal{G}$.

\medskip\noindent
For $\mathcal{F}$ in
$\textit{Mod}(\mathcal{X}_\etale, \mathcal{O}_\mathcal{X})$
there is a canonical map $\pi_F^*(\mathcal{F}|_{F_\etale}) \to \mathcal{F}$,
namely the map adjoint to the identification
$\mathcal{F}|_{F_\etale} = \pi_{F, *}\mathcal{F}$.
We will show that this map is an isomorphism if $\mathcal{F}$ is
a quasi-coherent module on $\mathcal{X}$.
Choose a scheme $U$ and a surjective \'etale morphism $U \to F$.
Denote $x : U \to \mathcal{X}$ the corresponding object of $\mathcal{X}$
over $U$. It suffices to show that
$\pi_F^*(\mathcal{F}|_{F_\etale}) \to \mathcal{F}$ is an isomorphism
after restricting to $\mathcal{X}_\etale/x = (\Sch/U)_\etale$.
Since $U \to F$ is \'etale, it follows from Remark \ref{remark-compare} that
$$
\pi_F^*(\mathcal{F}|_{F_\etale})|_{\mathcal{X}_\etale/x} =
\pi_U^*(\mathcal{F}|_{U_\etale})
$$
and that the restriction of the map
$\pi_F^*(\mathcal{F}|_{F_\etale}) \to \mathcal{F}$
to $\mathcal{X}_\etale/x = (\Sch/U)_\etale$
is equal to the corresponding map
$\pi_U^*(\mathcal{F}|_{U_\etale}) \to \mathcal{F}|_{(\Sch/U)_\etale}$.
Since we have seen the result is true for schemes in
Descent, Section \ref{descent-section-quasi-coherent-sheaves}\footnote{Namely,
if $U$ is a scheme and $\mathcal{F}$ is quasi-coherent on
$(\Sch/U)_\etale$, then $\mathcal{F} = \mathcal{H}^a$ for some
quasi-coherent module $\mathcal{H}$ on the scheme $U$ by
Descent, Proposition \ref{descent-proposition-equivalence-quasi-coherent}.
In other words, $\mathcal{F} = (\text{id}_{\etale,Zar})^*\mathcal{H}$
by Descent, Remark \ref{descent-remark-change-topologies-ringed-sites}
with notation as in Descent, Lemma \ref{descent-lemma-compare-sites}.
Then we have
$\text{id}_{\etale,Zar} = \pi_U \circ \text{id}_{small,\etale,Zar}$
and hence we see that $\mathcal{F} = \pi_U^*\mathcal{G}$ where
$\mathcal{G} = (\text{id}_{small,\etale,Zar})^*\mathcal{H}$ is
quasi-coherent. Then
$\pi_U^*i_U^*\mathcal{F} = \pi_U^*i_U^*\pi_U^*\mathcal{G} =
\pi_U^*\mathcal{G} = \mathcal{F}$ as desired.}
we conclude.

\medskip\noindent
Compatibility with pullbacks follows from the fact that
the quasi-inverse is given by $\pi_F^*$ and the
commutative diagram of ringed topoi in
Lemma \ref{lemma-compare-morphism}.
\end{proof}

\end{document}
